\chapter{Understanding and Preparing Data}
\label{ch:data}

\chapterguide%
  {\chapsref{ch:orientation}{ch:python}: the workflow, computational vocabulary, and the distinction between features and labels.}
  {clean data, choose valid splits, handle missingness and categories, engineer features, and prevent leakage.}

Reliable learning starts with trustworthy data: clear labels, valid measurements, and a split that reflects future use.

Ask: \emph{What does the data mean? What is wrong? What must be learned?}

\section{Start with a data quality assessment}

\begin{mlsteps}
  \item \term{What does one row represent?} A person, transaction, image, day, machine, or something else?
  \item \term{What does each column mean?} Record its type, unit, valid range, and source.
  \item \term{When is prediction made?} Features must exist then. Lab~1's G1 and G2 grades are available at year-end, not enrolment.
  \item \term{Which rows are related?} Repeated measurements from one person or device must usually stay in the same split.
  \item \term{Who or what is missing?} The dataset may not represent the full population where the model will be used.
\end{mlsteps}

\section{A safe order of operations}

\begin{mlfigure}
\begin{tikzpicture}[
  node distance=0.45cm,
  box/.style={draw=MLBlue,fill=MLPaleBlue,rounded corners=2mm,minimum width=2.75cm,minimum height=1.05cm,align=center,font=\scriptsize\bfseries\color{MLNavy}},
  arr/.style={-{Stealth[length=2mm]},thick,draw=MLTeal}
]
\node[box] (meaning) {1. Define meaning\\rows, keys, units, timing};
\node[box,right=of meaning] (rules) {2. Apply fixed rules\\parse, standardize, validate};
\node[box,right=of rules] (split) {3. Split honestly\\groups, time, population};
\node[box,right=of split] (learn) {4. Learn preparation\\from training data only};
\draw[arr] (meaning)--(rules);
\draw[arr] (rules)--(split);
\draw[arr] (split)--(learn);
\end{tikzpicture}
\end{mlfigure}

\begin{mltable}
\begin{tabularx}{\linewidth}{@{}>{\bfseries\leavevmode\color{MLNavy}}p{0.18\linewidth}p{0.35\linewidth}X@{}}
\toprule
\rowcolor{MLPaleBlue!92}
Stage & Examples & Why it belongs there \\
\midrule
Before the split & Validate columns, keys, and values; parse dates; standardize units and whitespace. & These are fixed rules. They do not use the label or any statistic estimated from the data. \\
After the split & Learn imputation, scaling, category mappings, outlier thresholds, feature selection, and resampling rules. & Fit data-dependent choices on training rows only. \\
After modeling & Evaluate test performance once, document limits, and monitor data quality and drift. & Keep the test set for final evaluation, not design decisions. \\
\bottomrule
\end{tabularx}
\end{mltable}

\begin{keyidea}
Fixed rules, such as unit conversions, may precede splitting. Estimate medians, scales, category mappings, and cutoffs from training rows only.
\end{keyidea}

\section{Feature and label types}

\begin{mltable}
\begin{tabularx}{\linewidth}{@{}>{\bfseries\leavevmode\color{MLNavy}}p{0.20\linewidth}>{\raggedright\arraybackslash}p{0.27\linewidth}>{\raggedright\arraybackslash}X@{}}
\toprule
\rowcolor{MLPaleBlue!92}
Type & Example & Typical treatment \cr
\midrule
Numerical & age, temperature, income & keep numeric; scale or transform when needed \cr
Nominal category & country, product type & one-hot encode or use a model that handles categories \cr
Ordinal category & low, medium, high & encode in the true order, but do not assume equal gaps without reason \cr
Date or time & purchase date, hour & extract useful parts such as elapsed time, weekday, or season \cr
Text & review, message & convert to numerical counts or other classical text features \cr
\bottomrule
\end{tabularx}
\end{mltable}

Labels usually fall into two basic groups:

\begin{itemize}
  \item A \term{continuous number} suggests regression, such as price or temperature.
  \item A \term{category} suggests classification, such as fraud/not fraud or animal species.
\end{itemize}

\section{Clean structural problems before modeling}

Make stored values match their documented meaning using reproducible cleaning code.

\subsection{Confirm row identity and handle duplicates}

\begin{mltable}
\renewcommand{\arraystretch}{1.18}
\setlength{\tabcolsep}{5pt}
\begin{tabularx}{\linewidth}{@{}>{\footnotesize\bfseries\leavevmode\color{MLNavy}}p{0.20\linewidth}>{\footnotesize}p{0.35\linewidth}>{\footnotesize}X@{}}
\toprule
\rowcolor{MLPaleBlue!92}
\normalfont\bfseries Pattern & \bfseries What it may mean & \bfseries Safe response \\
\midrule
Exact repeated row & Duplicate import or distinct events lacking identifiers & Remove only if repetition is meaningless; log the count. \\
Repeated key & Key conflict, multiple versions, or valid one-to-many data & Investigate: keep the latest, aggregate, repair the key, or redefine rows. \\
Same person, different visits & Valid repeated measurements & Keep rows; group splits by person to prevent overlap. \\
Near duplicate & Spelling, rounding, or record-linkage differences & Define matching rules and review uncertain matches. \\
\bottomrule
\end{tabularx}
\end{mltable}

\subsection{Repair types, labels, dates, and units}

\begin{mlsteps}
  \item \term{Preserve the raw value.} Keep the source file unchanged and create a cleaned copy or scripted transformation.
  \item \term{Parse explicitly.} Convert numbers and dates with known rules, and count values that fail to parse.
  \item \term{Normalize representation.} Trim accidental whitespace, standardize case where case has no meaning, and map documented synonyms to one label.
  \item \term{Normalize units.} Convert to one stated unit using a known conversion; never mix unit conversion with statistical scaling.
  \item \term{Validate the result.} Check allowed categories, ranges, key uniqueness, and date order with assertions.
\end{mlsteps}

\subsection{Distinguish impossible values from unusual values}

\term{Impossible values}, such as negative ages, violate domain rules. \term{Rare values}, such as large purchases, may be valid.

\begin{mltable}
\begin{tabularx}{\linewidth}{@{}>{\bfseries\leavevmode\color{MLNavy}}p{0.15\linewidth}>{\raggedright\arraybackslash}X@{}}
\toprule
\rowcolor{MLPaleBlue!92}
Action & When to use it \\
\midrule
Correct & Only when an authoritative source or deterministic rule reveals the true value. \\
Set missing & When the stored value is known to be invalid but the true value is unknown. Impute later inside the training pipeline if appropriate. \\
Drop & When the row cannot answer the project question and a documented inclusion rule justifies removal. Report the count and check whether exclusions are concentrated in a group. \\
Flag & When a suspicious value may still be real. A binary quality flag lets the model or analyst retain that information. \\
Keep & When a value is valid, even if it is inconvenient. Rare cases may be exactly where a useful model is most needed. \\
\bottomrule
\end{tabularx}
\end{mltable}

\begin{pitfall}
Do not remove difficult cases merely to improve scores; they may occur after deployment.
\end{pitfall}

\subsection{Keep a cleaning log}

Log each rule's reason, columns, affected row count, and action: correction, missingness, exclusion, or flagging. Retain identifiers for tracing changes, not as model features.

\section{Training, validation, and test sets}
\label{sec:data-splits}

\begin{mlfigure}
\begin{tikzpicture}[x=0.093\linewidth,font=\small]
  \draw[fill=MLBlue!28,draw=MLBlue] (0,0) rectangle (6.5,0.85);
  \draw[fill=MLTeal!22,draw=MLTeal] (6.5,0) rectangle (8.3,0.85);
  \draw[fill=MLOrange!22,draw=MLOrange] (8.3,0) rectangle (10,0.85);
  \node[align=center] at (3.25,0.425) {Training\\learn parameters};
  \node[align=center,font=\footnotesize] at (7.4,0.425) {Validation\\choose settings};
  \node[align=center,font=\footnotesize] at (9.15,0.425) {Test\\final check};
\end{tikzpicture}
\end{mlfigure}

\begin{description}[leftmargin=1.15in,style=nextline]
  \item[Training set] Used to fit model parameters and data-dependent preprocessing steps.
  \item[Validation set] Used to compare models, thresholds, and hyperparameters.
  \item[Test set] Used only for the final estimate after all choices are complete.
\end{description}

Training often uses 60--80\% of rows. Choose proportions for dataset size while preserving separate validation and test roles.

\section{Choose a split that imitates the future}

A random split is useful only when future rows are independent and drawn in the same way.

\begin{description}[leftmargin=1.35in,style=nextline]
  \item[Stratified split] Keeps class proportions similar across subsets.
  \item[Group split] Keeps all rows from one person, device, school, or customer in one subset.
  \item[Time split] Trains on the past and evaluates on later periods.
  \item[Spatial split] Separates locations when nearby observations are strongly related.
\end{description}

\begin{workedexample}
With five visits per patient, a random row split may let the model recognize patients seen in training. Group by patient when evaluating predictions for new patients.
\end{workedexample}

\begin{pitfall}
An unrealistic split gives an unrealistic performance estimate.
\end{pitfall}

\section{Data leakage}
\label{sec:data-leakage}

\term{Data leakage} happens when information the model would not have at prediction time, or information from the held-out rows, finds its way into the features, the preprocessing, the choice of model, or the split.

Examples include:

\begin{itemize}
  \item predicting hospital readmission using a field recorded after discharge;
  \item calculating the mean for standardization using the entire dataset;
  \item selecting features using test-set performance;
  \item placing near-duplicate rows in both training and test data;
  \item using future values to predict an earlier point in time;
  \item including an identifier or proxy that effectively reveals the label, such as a claim-resolution code when predicting whether a claim will be approved.
\end{itemize}

The safe rule is simple: fit every data-dependent choice on training data only, then apply the learned transformation unchanged to validation and test data.

The diagram contrasts fitting the preprocessing before and after the split. Only the training rows should decide what a fitted transformation does.

\begin{mlfigure}
\begin{tikzpicture}[
  node distance=0.35cm and 0.5cm,
  ok/.style={draw=MLGreen,fill=MLPaleTeal,rounded corners=1.5mm,minimum width=2.85cm,minimum height=0.78cm,inner xsep=3pt,align=center,font=\scriptsize\bfseries\color{MLNavy}},
  bad/.style={draw=MLRed,fill=MLPaleRed,rounded corners=1.5mm,minimum width=2.85cm,minimum height=0.78cm,inner xsep=3pt,align=center,font=\scriptsize\bfseries\color{MLNavy}},
  arr/.style={-{Stealth[length=1.9mm]},thick,draw=MLTeal},
  barr/.style={-{Stealth[length=1.9mm]},thick,draw=MLRed},
  tag/.style={font=\scriptsize\bfseries,align=left}
]
\node[tag,color=MLRed] (t1) {WRONG};
\node[bad,right=0.35cm of t1] (w1) {Load all data};
\node[bad,right=of w1] (w2) {Fit prep on all data};
\node[bad,right=of w2] (w3) {Split};
\node[bad,right=of w3] (w4) {Train, test, celebrate};
\draw[barr] (w1)--(w2); \draw[barr] (w2)--(w3); \draw[barr] (w3)--(w4);
\node[font=\scriptsize\itshape\color{MLRed},below=0.12cm of w2] {test rows influenced the scaler};

\node[tag,color=MLGreen,below=1.35cm of t1] (t2) {RIGHT};
\node[ok,right=0.35cm of t2] (r1) {Load all data};
\node[ok,right=of r1] (r2) {Split first};
\node[ok,right=of r2] (r3) {Fit prep on training only};
\node[ok,right=of r3] (r4) {Transform test, then evaluate};
\draw[arr] (r1)--(r2); \draw[arr] (r2)--(r3); \draw[arr] (r3)--(r4);
\end{tikzpicture}
\end{mlfigure}

\section{Missing values: meaning and strategy}

Missingness may mean unmeasured, inapplicable, costly to collect, or related to the outcome.

\begin{mlsteps}
  \item Measure how much is missing in each feature and group.
  \item Investigate why it is missing and whether the pattern carries information.
  \item Choose a treatment and fit it on training rows only (\secref{sec:data-leakage}).
  \item Add a missing-value indicator when the fact that a value is missing may matter.
  \item Check whether the treatment changes important distributions.
\end{mlsteps}

Use numerical medians, a categorical missing value, or model-based imputation when justified.

\term{MCAR}: missing independently of values. \term{MAR}: missingness explained by observed values. \term{MNAR}: missingness depends on unobserved values. The table alone cannot identify the mechanism; investigate causes and compare treatments.

\section{Encoding categorical features}

Most algorithms need numerical inputs. The encoding must preserve the meaning of the category.

\begin{mltable}
\begin{tabularx}{\linewidth}{@{}>{\bfseries\leavevmode\color{MLNavy}}p{0.35\linewidth}X@{}}
\toprule
\rowcolor{MLPaleBlue!92}
Encoding method & How it works \\
\midrule
One-hot encoding & Creates one binary column per category. It is a safe default for nominal features with a manageable number of values. \\
Ordinal encoding & Uses ordered numbers only when the categories truly have an order. \\
Frequency encoding & Replaces a category with how often it appears in training data. \\
Label-informed mean encoding & Uses label statistics; apply the leakage rule from \secref{sec:data-leakage}. \\
\bottomrule
\end{tabularx}
\end{mltable}

Define how the pipeline handles unseen categories, such as an unknown-category value.

\section{Scaling and transforming numerical features}
\label{sec:scaling}

Scaling changes a feature's measuring stick. Distances, margins, penalties, and gradient steps otherwise let large numerical units dominate smaller ones.

The three course scalers answer three different questions:

\begin{mltable}
\begin{tabularx}{\linewidth}{@{}>{\bfseries\leavevmode\color{MLNavy}}p{0.20\linewidth}p{0.29\linewidth}X@{}}
\toprule
\rowcolor{MLPaleBlue!92}
Method & Transformation & Best first use \\
\midrule
Standard scaling & $z=(x-\mu)/\sigma$ & General default for SVM, $k$-NN, clustering, and regularized linear models when mean and standard deviation are meaningful. \\
Min--max scaling & $x'=(x-x_{\min})/(x_{\max}-x_{\min})$ & Put a feature into $[0,1]$ when a bounded range is convenient and extreme values are not dominating the observed range. \\
Robust scaling & $x'=(x-\operatorname{median}(x))/\operatorname{IQR}(x)$ & Data with strong outliers, because the median and interquartile range are less affected by extreme observations. \\
\bottomrule
\end{tabularx}
\end{mltable}

Fit every scaler on training rows only (\secref{sec:data-leakage}).

\begin{workedexample}
For $[2400,36,23075,70,1,905]$, min--max scaling maps 1 and 23,075 to 0 and 1. Robust scaling uses the median and IQR, reducing sensitivity to the extreme value.
\end{workedexample}

\begin{mltable}
\begin{tabularx}{\linewidth}{@{}>{\bfseries\leavevmode\color{MLNavy}}p{0.27\linewidth}p{0.20\linewidth}X@{}}
\toprule
\rowcolor{MLPaleBlue!92}
Model family & Scale-sensitive? & Why \\
\midrule
$k$-NN, $K$-means, DBSCAN & Usually yes & Their neighborhoods or clusters are defined by distance. \\
SVM / SVR & Usually yes & Margin and kernel distances depend directly on feature magnitude. \\
Ridge, lasso, elastic net, logistic regression & Usually yes & The penalty compares coefficient sizes across features; unequal units distort that comparison. \\
Decision trees and tree ensembles & Usually no & Splits depend on ordering and thresholds, which monotonic rescaling preserves. \\
Naive Bayes & Depends on variant & Gaussian naive Bayes does not need it; other variants expect counts or binary inputs, so match the preprocessing to the variant. \\
\bottomrule
\end{tabularx}
\end{mltable}

A $\log(1+x)$ or power transform changes distribution shape; scaling mainly changes units and location.

\begin{intuition}
Income values of 20,000--300,000 can dominate ages of 10--80 in Euclidean distance. Scaling makes their numerical ranges comparable; it does not establish their relevance.
\end{intuition}

% === VISUAL ADDITION: scaling-geometry ===
\subsection{Scaling changes the geometry, drawn}

Scaling changes the numerical gaps used by distance-based models.

\begin{mlfigure}
\begin{tikzpicture}[font=\scriptsize]
\begin{scope}
  \fill[MLPaleRed!30,rounded corners=1mm] (0,0) rectangle (6.25,3.45);
  \draw[rounded corners=1mm,draw=MLRed!45] (0,0) rectangle (6.25,3.45);
  \node[font=\small\bfseries\color{MLNavy}] at (3.125,3.05) {Before scaling};

  \node[anchor=west,text=MLMuted] at (0.35,2.30) {age gap};
  \draw[MLBlue,line width=5pt] (2.05,2.30) -- (2.85,2.30);
  \node[anchor=west,text=MLBlue] at (3.07,2.30) {$10$ years};

  \node[anchor=west,text=MLMuted] at (0.35,1.48) {income gap};
  \draw[MLOrange,line width=5pt] (2.05,1.48) -- (5.50,1.48);
  \node[anchor=east,text=MLOrange] at (5.75,1.80) {$40{,}000$ dollars};

  \node[align=center,text width=5.35cm,text=MLRed] at (3.125,0.55)
    {Raw Euclidean distance is almost entirely the income column.};
\end{scope}

\begin{scope}[xshift=6.75cm]
  \fill[MLPaleTeal!45,rounded corners=1mm] (0,0) rectangle (6.25,3.45);
  \draw[rounded corners=1mm,draw=MLTeal!55] (0,0) rectangle (6.25,3.45);
  \node[font=\small\bfseries\color{MLNavy}] at (3.125,3.05) {After standardization};

  \node[anchor=west,text=MLMuted] at (0.35,2.30) {age gap};
  \draw[MLBlue,line width=5pt] (2.05,2.30) -- (4.50,2.30);
  \node[anchor=west,text=MLBlue] at (4.72,2.30) {$0.50$ SD};

  \node[anchor=west,text=MLMuted] at (0.35,1.48) {income gap};
  \draw[MLTeal,line width=5pt] (2.05,1.48) -- (4.01,1.48);
  \node[anchor=west,text=MLTeal] at (4.23,1.48) {$0.40$ SD};

  \node[align=center,text width=5.35cm,text=MLGreen] at (3.125,0.55)
    {Both features now contribute on comparable measuring sticks.};
\end{scope}

\node[align=center,text width=12.5cm,font=\scriptsize\bfseries\color{MLNavy}] at (6.50,-0.48)
  {Same observations; different geometry. Scaling changes distance, not the ordering within a feature.};
\end{tikzpicture}
\end{mlfigure}

\begin{keyidea}
Changing units can change neighbors, margins, and clusters without changing observations or labels.
\end{keyidea}
% === END VISUAL ADDITION ===

\section{Outliers, imbalance, and rare cases}

Outliers may be valid rare events or evidence of mixed populations. Investigate before deleting them; learn percentile or IQR cutoffs from training rows only.

Record class counts and retain rare cases. Stratify when random splitting is appropriate.

\begin{keyidea}
Choose evaluation metrics before an imbalance intervention (\chapref{ch:evaluation}). \chapref{ch:training} covers weights, thresholds, and resampling, including SMOTE; \chapref{ch:knn-nb} explains its neighbor geometry.
\end{keyidea}

\section{Feature engineering}

\term{Feature engineering} creates useful inputs from raw data. Examples include:

\begin{itemize}
  \item price per unit, body-mass index, or another meaningful ratio;
  \item elapsed time between two dates;
  \item weekday, month, or cyclical time features;
  \item interactions such as area multiplied by location score;
  \item counts or averages from past behavior, calculated without using the future.
\end{itemize}

\subsection*{Feature engineering by data type}
Choose transformations that preserve meaning and use only prediction-time information.

\begin{mltable}
\begin{tabularx}{\linewidth}{@{}>{\bfseries\leavevmode\color{MLNavy}}p{0.19\linewidth}p{0.31\linewidth}X@{}}
\toprule
\rowcolor{MLPaleBlue!92}
Data type & Common transformation ideas & Main safety question \\
\midrule
Categorical & One-hot or ordinal encoding; rare-category grouping & Does encoding respect category order and use training-only mappings? \\
Numeric & Logs, powers, ratios, interactions, polynomials, scaling & Are units and meaning respected, with parameters fitted on training data only? \\
Temporal & Lags, elapsed time, rolling summaries, seasonality & Were all inputs available at prediction time? \\
Text or images & Counts, embeddings, handcrafted or learned features & Is task-relevant structure retained without label or future leakage? \\
\bottomrule
\end{tabularx}
\end{mltable}

\subsection{A safety test for engineered features}
For each derived feature, ask what information it uses and when that information becomes available.

\begin{mltable}
\begin{tabularx}{\linewidth}{@{}>{\bfseries\leavevmode\color{MLNavy}}p{0.26\linewidth}p{0.30\linewidth}X@{}}
\toprule
\rowcolor{MLPaleBlue!92}
Feature type & Example & Safe treatment \\
\midrule
Deterministic & BMI from height and weight; hour from timestamp & Compute reproducibly from legitimately available inputs. \\
Distribution-fitted & Power transforms, frequency encoding, clipping, category grouping & Fit on training rows only, usually in a pipeline. \\
Label-aware & Mean encoding, supervised selection, label-based category merging & Fit within training folds, never on all labels before validation (\chapref{ch:training}). \\
Temporal aggregate & Lags, rolling means, prior-week counts & Use only data available at prediction time; never fill past gaps with future values. \\
\bottomrule
\end{tabularx}
\end{mltable}

\section{Use a preprocessing pipeline}
\label{sec:preprocessing-pipeline}

A pipeline fits transformations on training rows and reuses them on new data. Keep feature selection and resampling within training folds too (\chapref{ch:training}).

\begin{mlfigure}
\begin{tikzpicture}[
  node distance=0.35cm,
  box/.style={draw=MLBlue,fill=MLPaleBlue,rounded corners=2mm,minimum height=0.76cm,inner xsep=5pt,align=center,font=\scriptsize\color{MLNavy}},
  arr/.style={-{Stealth[length=2mm]},thick,draw=MLTeal}
]
\node[box,minimum width=1.8cm] (raw) {Raw data};
\node[box,right=of raw,minimum width=3.5cm] (check) {Fixed schema and quality rules};
\node[box,right=of check,minimum width=2.1cm] (split) {Honest split};
\node[box,right=of split,minimum width=3.4cm] (prep) {Train-fitted preprocessing};
\node[box,right=of prep,minimum width=1.8cm] (model) {Fit model};
\draw[arr] (raw)--(check); \draw[arr] (check)--(split); \draw[arr] (split)--(prep); \draw[arr] (prep)--(model);
\end{tikzpicture}
\end{mlfigure}

Document the data source, collection period, row definition, feature units, label timing, inclusion rules, split logic, random seed, known bias, and privacy limits.

\begin{keyidea}
Validate the whole pipeline, refitting it inside each cross-validation training fold (\chapref{ch:evaluation}).
\end{keyidea}

\begin{debugtask}
A teammate computes median values, standardization statistics, and a label-informed mean encoding on the complete labeled table, then creates the train--test split. Identify every quantity that carried information from the future holdout backward. Rewrite the order of operations in six lines of pseudocode.
\end{debugtask}

\labsection{1}{Feature engineering and feature selection}{sec:lab01-features}

\begin{labproject}{Check the tools, understand the data, build features, and choose among them}
\textbf{Goal.} Check the environment, engineer features, compare scaling, and select features without leakage.

\medskip\noindent\textbf{Setup.}\par
\begin{lstlisting}[style=mlpython]
from pathlib import Path

import matplotlib
import matplotlib.pyplot as plt
import numpy as np
import pandas as pd
import sklearn
from sklearn.compose import ColumnTransformer
from sklearn.feature_selection import SelectKBest, f_regression
from sklearn.impute import SimpleImputer
from sklearn.linear_model import Lasso, LogisticRegression, Ridge
from sklearn.metrics import mean_absolute_error, pairwise_distances, roc_auc_score
from sklearn.model_selection import train_test_split
from sklearn.pipeline import Pipeline
from sklearn.preprocessing import OneHotEncoder, StandardScaler

# Notebooks live in all_experiments/, so the course folder is one level up.
PROJECT_ROOT = Path.cwd().parent
DATA = PROJECT_ROOT / "all_datasets"
\end{lstlisting}

\medskip\noindent\textbf{Part A --- Environment check and a first look at the data.}\par

Record package versions for reproducibility.

\begin{lstlisting}[style=mlpython]
print("NumPy:", np.__version__)
print("pandas:", pd.__version__)
print("Matplotlib:", matplotlib.__version__)
print("scikit-learn:", sklearn.__version__)
\end{lstlisting}

Titanic has one row per passenger; \texttt{Survived}=1 denotes survival. Inspect shape, types, missingness, and class balance.

\begin{lstlisting}[style=mlpython]
df = pd.read_csv(DATA / "titanic_dataset.csv")

print("shape:", df.shape)
print(df.head())
print("\nData types:\n", df.dtypes)
print("\nSurvival rate:\n", df["Survived"].value_counts(normalize=True))
\end{lstlisting}

Plot missing-value percentages by column.

\begin{lstlisting}[style=mlpython]
missing = df.isna().mean().sort_values(ascending=False)
missing = missing[missing > 0]

plt.figure(figsize=(5, 3))
plt.barh(missing.index, missing.values * 100)
plt.xlabel("Share of rows missing (%)")
plt.title("Titanic: where the gaps are")
plt.tight_layout()
plt.show()
\end{lstlisting}

\medskip\noindent\textbf{Part B --- Feature engineering.}\par

Derive family size, isolation, title, and log fare from prediction-time information. Each feature must exist before the outcome.

\begin{lstlisting}[style=mlpython]
df["FamilySize"] = df["SibSp"] + df["Parch"] + 1
df["IsAlone"] = (df["FamilySize"] == 1).astype(int)
df["Title"] = df["Name"].str.extract(r",\s*([^\.]+)\.")[0].str.strip()
df["Title"] = df["Title"].where(df["Title"].isin(["Mr", "Mrs", "Miss", "Master"]), "Other")
df["LogFare"] = np.log1p(df["Fare"])

print(df[["Name", "Title", "FamilySize", "IsAlone", "Fare", "LogFare"]].head())
print("\nTitle counts:\n", df["Title"].value_counts())
\end{lstlisting}

Compare survival rates across engineered features.

\begin{lstlisting}[style=mlpython]
fig, axes = plt.subplots(1, 2, figsize=(9.5, 3.4))
df.groupby("Title")["Survived"].mean().sort_values().plot.barh(ax=axes[0])
axes[0].set_xlabel("Survival rate")
axes[0].set_title("Title carries the sex-and-age story in one column")
df.groupby("FamilySize")["Survived"].mean().plot.bar(ax=axes[1], rot=0)
axes[1].set_ylabel("Survival rate")
axes[1].set_title("Small families did best; large ones did worst")
plt.tight_layout()
plt.show()
\end{lstlisting}

\textbf{Inspect.} Survival varies by title and family size; small families fare better than lone passengers or large families.

Compare original and engineered features using the same model and split. Fit imputation, scaling, and encoding inside the training pipeline.

\begin{lstlisting}[style=mlpython]
def make_preprocess(num_cols, cat_cols):
    return ColumnTransformer([
        ("num", Pipeline([
            ("impute", SimpleImputer(strategy="median")),
            ("scale", StandardScaler()),
        ]), num_cols),
        ("cat", Pipeline([
            ("impute", SimpleImputer(strategy="most_frequent")),
            ("onehot", OneHotEncoder(handle_unknown="ignore")),
        ]), cat_cols),
    ])

y = df["Survived"]
feature_sets = {
    "raw columns": (["Age", "SibSp", "Parch", "Fare"], ["Pclass", "Sex", "Embarked"]),
    "raw + engineered": (["Age", "SibSp", "Parch", "LogFare", "FamilySize", "IsAlone"],
                         ["Pclass", "Sex", "Embarked", "Title"]),
}

for name, (num_cols, cat_cols) in feature_sets.items():
    X = df[num_cols + cat_cols]
    X_train, X_valid, y_train, y_valid = train_test_split(
        X, y, test_size=0.25, random_state=42, stratify=y
    )
    model = Pipeline([
        ("preprocess", make_preprocess(num_cols, cat_cols)),
        ("model", LogisticRegression(max_iter=2000)),
    ]).fit(X_train, y_train)
    auc = roc_auc_score(y_valid, model.predict_proba(X_valid)[:, 1])
    print(f"{name:18s} validation ROC AUC = {auc:.3f}")
\end{lstlisting}

\textbf{Inspect.} Higher validation AUC measures the gain from features, not a different estimator.

\medskip\noindent\textbf{Part C --- Scaling changes the geometry.}\par

Breast-cancer measurements have unequal scales: area reaches thousands while smoothness is below one. Compare their effect on distances.

\begin{lstlisting}[style=mlpython]
cancer = pd.read_csv(DATA / "breast_cancer_dataset.csv")
feature_cols = [c for c in cancer.columns if c not in {"id", "diagnosis", "Unnamed: 32"}]
X_cancer = cancer[feature_cols]

three = ["area_mean", "texture_mean", "smoothness_mean"]
print(X_cancer[three].describe().loc[["mean", "std"]].round(3))
\end{lstlisting}

Compare raw and standardized distances among the first five tumors; darker heatmap cells indicate closer pairs.

\begin{lstlisting}[style=mlpython]
raw = pairwise_distances(X_cancer.iloc[:5], metric="euclidean")
X_scaled = StandardScaler().fit_transform(X_cancer)
scaled = pairwise_distances(X_scaled[:5], metric="euclidean")

fig, axes = plt.subplots(1, 2, figsize=(9, 3.6))
for ax, matrix, title in [(axes[0], raw, "Raw units"), (axes[1], scaled, "Standardized")]:
    im = ax.imshow(matrix, cmap="viridis")
    ax.set_title(title)
    ax.set_xlabel("tumour")
    ax.set_ylabel("tumour")
    for i in range(5):
        for j in range(5):
            ax.text(j, i, f"{matrix[i, j]:.0f}", ha="center", va="center", color="white", fontsize=8)
    fig.colorbar(im, ax=ax, shrink=0.8)
plt.suptitle("Same five tumours, two geometries")
plt.tight_layout()
plt.show()
\end{lstlisting}

\textbf{Inspect.} The closest pair changes from tumors 3/5 to 2/5 after scaling. Units can change neighbors.

\medskip\noindent\textbf{Part D --- Feature selection without leakage.}\par

Predict \texttt{G3} for 649 students at enrolment. Exclude later grades \texttt{G1} and \texttt{G2}, leaving 30 candidate columns.

\begin{lstlisting}[style=mlpython]
students = pd.read_csv(DATA / "student_performance_dataset.csv")
print("shape:", students.shape, " missing:", students.isna().sum().sum(),
      " duplicates:", students.duplicated().sum())

target = "G3"
leaky = ["G1", "G2"]
X_s = students.drop(columns=[target] + leaky)
y_s = students[target]
s_num = X_s.select_dtypes(include="number").columns.tolist()
s_cat = X_s.select_dtypes(exclude="number").columns.tolist()
print(len(s_num), "numeric and", len(s_cat), "text features available at enrolment")

Xs_train, Xs_valid, ys_train, ys_valid = train_test_split(X_s, y_s, test_size=0.25, random_state=42)
\end{lstlisting}

Compare ridge using all features or ten \texttt{SelectKBest} features, and lasso's sparse fit. Keep selection inside training pipelines. A fourth run deliberately adds \texttt{G2} to expose leakage.

\begin{lstlisting}[style=mlpython]
def evaluate(pipeline, X_tr, X_va, label):
    pipeline.fit(X_tr, ys_train)
    mae = mean_absolute_error(ys_valid, pipeline.predict(X_va))
    print(f"{label:32s} validation MAE = {mae:.3f} grade points")
    return mae

results = {}
results["all 30 features"] = evaluate(Pipeline([
    ("preprocess", make_preprocess(s_num, s_cat)),
    ("model", Ridge(alpha=1.0)),
]), Xs_train, Xs_valid, "all 30 features")

results["SelectKBest, k=10"] = evaluate(Pipeline([
    ("preprocess", make_preprocess(s_num, s_cat)),
    ("select", SelectKBest(f_regression, k=10)),
    ("model", Ridge(alpha=1.0)),
]), Xs_train, Xs_valid, "SelectKBest, k=10")

lasso = Pipeline([
    ("preprocess", make_preprocess(s_num, s_cat)),
    ("model", Lasso(alpha=0.1)),
])
results["lasso (selects itself)"] = evaluate(lasso, Xs_train, Xs_valid, "lasso (selects itself)")

leak_num = s_num + ["G2"]
results["all 30 + G2 (leakage)"] = evaluate(Pipeline([
    ("preprocess", make_preprocess(leak_num, s_cat)),
    ("model", Ridge(alpha=1.0)),
]), Xs_train.assign(G2=students.loc[Xs_train.index, "G2"]),
    Xs_valid.assign(G2=students.loc[Xs_valid.index, "G2"]), "all 30 + G2 (leakage)")
\end{lstlisting}

\begin{lstlisting}[style=mlpython]
coefs = pd.Series(lasso.named_steps["model"].coef_,
                  index=lasso.named_steps["preprocess"].get_feature_names_out())
kept = coefs[coefs != 0].sort_values()
print(f"lasso kept {len(kept)} of {len(coefs)} encoded columns")

fig, axes = plt.subplots(1, 2, figsize=(11, 3.8))
axes[0].barh(list(results), list(results.values()))
axes[0].set_xlabel("Validation MAE (grade points)")
axes[0].set_title("Fewer features, same error; a leaked feature, half the error")
axes[1].barh(kept.index, kept.values)
axes[1].axvline(0, color="grey", linewidth=0.8)
axes[1].set_title("What the lasso kept")
plt.tight_layout()
plt.show()
\end{lstlisting}

\textbf{Inspect.} Selected features give similar or better error with fewer columns. Adding \texttt{G2} cuts error sharply but invalidates an enrolment-time forecast. Inspect lasso's retained school, failure, education, alcohol, and study-time features without treating them as causal.

\begin{tryit}
Add cabin initial or fare per family member; compare Part~B validation AUC. Rerun Part~D with \texttt{k=5} and \texttt{k=20}; compare errors.
\end{tryit}
\end{labproject}


\begin{quickcheck}
\begin{enumerate}
  \item Why must train/test splitting happen before any preprocessing step that learns from the data?
  \item What is the difference between an identifier, a feature, and a label?
  \item Why do you need to know the moment of prediction before you can say whether a feature is leakage?
\end{enumerate}
\end{quickcheck}
